\documentclass[12pt]{article}
\usepackage[utf8]{inputenc}
\usepackage[T1]{fontenc}
\usepackage{hyperref}
\title{Wolani Language Resources for K-12 Curriculum Development}
\author{Compiled by Claude Sonnet 5.5 (claude-sonnet-5-5) for Andr\'as Kornai}
\date{}
\begin{document}
\maketitle

\section{Introduction}
Wolani (also Wodani; ISO 639-3 wod, Glottocode wola1243) is a Papuan language of the Paniai lakes region, Central Papua, Indonesia, with about 5{,}000 speakers. Wikipedia says it is related to Moni, Ekari, Auye and Dao and that documentation is limited. Do not confuse it with the Moni (Migani) language. Search found almost nothing usable for schools.

\section{Primary Educational Resources}
None found. The language is not known to be taught in schools.

\section{Linguistic and Reference Materials}
Wikipedia has a language article: \url{https://en.wikipedia.org/wiki/Wolani_language}. It also has an article on the Wolani people: \url{https://en.wikipedia.org/wiki/Wolani_people}.

PARADISEC holds an Arthur Capell manuscript of Wolani and Western Dani sentences and phrases: \url{https://catalog.paradisec.org.au/repository/AC2/IJWDAN104}.

\section{Religious and Community Resources}
None found.

\section{Digital and Audio Resources}
No audio or software resources were found, apart from the PARADISEC catalogue item above.

\section{Curriculum Development Context}
Documentation of Wolani is very thin. A curriculum would need basic descriptive work, an orthography and a primer, done with the community and Indonesian education authorities in Papua.
\end{document}